\section{Clinical Deployment Assessment}
\label{sec:deployment}

Deploying a medical imaging classifier in healthcare is governed not by benchmark AUC alone, but by clinical workflow fit, patient safety, regulatory compliance, privacy law, demographic fairness, clinician trust, and lifecycle engineering~\cite{Esteva2019guide,Finlayson2021clinician}. Below we present a structured feasibility assessment for deploying our adapted dermatoscopic triage system in real clinical practice.

\subsection{Clinical Problem, Stakeholders, and Workflow Integration}
The target clinical problem is early triage of pigmented and non-pigmented skin lesions to prioritize urgent dermatology referral and biopsy for malignant lesions (melanoma, basal cell carcinoma, and actinic keratosis/intraepithelial carcinoma) while curbing unnecessary excisions of benign keratoses and melanocytic nevi. Three primary stakeholder groups interact with the system: (1)~\emph{primary-care physicians and nurse practitioners}, who perform first-line skin checks without specialist dermoscopy training and face the highest risk of missed early melanomas; (2)~\emph{board-certified dermatologists and teledermatology readers}, who adjudicate referrals and decide biopsy margins; and (3)~\emph{patients and hospital clinical-safety officers}, who bear the consequences of delayed cancer diagnosis, biopsy morbidity, and healthcare cost.

Within the clinical workflow (Figure~\ref{fig:clinical_workflow}), the model is designed strictly as an \textbf{assistive triage decision-support tool}—never an autonomous diagnostic replacement. During a primary-care visit, a clinician captures a standardized polarized dermatoscopic image using a calibrated handheld dermatoscope integrated into the Electronic Health Record (EHR) or PACS via DICOM Visible Light Photography. The inference service returns three outputs within $<500$~ms: calibrated malignancy probability, operating-point triage flag (set at the high-sensitivity screening threshold $\tau_{\text{scr}}$), and side-by-side Grad-CAM / Attention Rollout overlays. Crucially, the clinician must document an independent clinical impression before viewing AI output—mitigating anchoring and automation bias—and retains unconditional override authority whenever clinical history (rapid lesion evolution, bleeding, pain, or ``ugly duckling'' discordance) warrants biopsy despite a low algorithmic score~\cite{Tschandl2020human}.

\subsection{Clinical Performance Significance Against Dermatologists and Benchmarks}
Evaluating clinical significance requires comparing our operating points against peer-reviewed human reader benchmarks rather than arbitrary machine learning thresholds. In the international diagnostic study by Tschandl et al.~\cite{Tschandl2019comparison} involving 511 human readers (including 283 board-certified dermatologists), human experts achieved 62.6\% specific 7-class accuracy (18.78/30) and roughly 89\% sensitivity at 75\% specificity for malignancy triage on static dermoscopic images. Similarly, Brinker et al.~\cite{Brinker2019deep} showed that deep convolutional classifiers outperformed 136 of 157 German dermatologists on dermoscopic melanoma classification, and Esteva et al.~\cite{Esteva2017dermatologist} demonstrated dermatologist-level ROC performance on biopsy-proven skin cancers.

On our untouched HAM10000 test set, \WSixBestModel{} achieves \WSixBestBalAcc{} 7-class balanced accuracy, \WSixBestMalAUC{} malignancy AUC-ROC, and \WSixBestSensScr{} sensitivity at \WSixBestSpecScr{} specificity under the validation-selected screening threshold—demonstrating competitive discrimination and positive net benefit on Decision Curve Analysis~\cite{Vickers2006decision} across practical biopsy thresholds ($p_t \in [0.05, 0.45]$). However, three critical caveats prevent claiming immediate clinical readiness. First, HAM10000 represents retrospective, largely European and Australian clinic cohorts (Vienna and Queensland); Tschandl et al.~\cite{Tschandl2019comparison} demonstrated that algorithmic accuracy drops far more steeply than dermatologist accuracy on out-of-distribution clinical centers. Second, FDA clearance for dermatology triage tools typically requires prospective, multi-center clinical reader studies establishing non-inferiority or superiority of \emph{clinician+AI} over \emph{clinician alone} on biopsy-confirmed endpoints, with predefined lower-bound sensitivity guarantees (commonly $\ge 90\%$--$95\%$ for invasive melanoma). Third, single-image classification omits patient longitudinal history and total-body context essential to real dermatologic judgment.

\subsection{Failure Mode Analysis and Patient Safety Controls}
Our qualitative and quantitative interpretability audit on the held-out test set (Section~\ref{sec:interpretability}, Table~\ref{tab:failure_modes}) documented five distinct failure modes with direct patient-safety implications:
\begin{enumerate}
  \item \textbf{Missed Melanoma (Hypopigmented or Atypical Mimic, FN):} Amelanotic or lightly pigmented melanomas lack classic dark reticular pigment networks and can be misclassified as benign nevi or keratoses. \emph{Patient harm:} Critical—delayed excision risks progression to stage~III/IV disease and mortality. \emph{Clinician catch:} Mandatory clinical history check (new/evolving lesion, bleeding) and polarized vascular pattern inspection override any negative AI screen.
  \item \textbf{Rare Non-Melanoma Malignancy False Negative (\texttt{akiec}/\texttt{bcc}, FN):} Superficial basal cell carcinomas or Bowen's disease with faint erythema can receive sub-threshold malignancy scores. \emph{Patient harm:} Moderate-to-High—local tissue destruction requiring Mohs reconstruction. \emph{Clinician catch:} Heightened suspicion on chronically sun-exposed head/neck sites in older adults ($>60$~years) and physical palpation for pearly borders or scaling.
  \item \textbf{Pigmented Benign Keratosis False Positive (\texttt{bkl} $\to$ \texttt{mel}, FP):} Heavily pigmented seborrheic keratoses trigger high malignancy probability due to dark color asymmetry. \emph{Patient harm:} Low-to-Moderate—patient anxiety, surgical scarring, and unnecessary dermatology workload. \emph{Clinician catch:} Recognition of pathognomonic milia-like cysts, comedo-like openings, and ``stuck-on'' sharp demarcation allows safe deferral of biopsy.
  \item \textbf{Artifact and Border Shortcut Attention:} Cases where Grad-CAM concentrates $>30\%$ energy on outer circular dermoscope rims, measurement rulers, or unremoved hair rather than lesion tissue. \emph{Patient harm:} High systemic risk—fragile predictions fail unpredictably when clinic hardware changes. \emph{Clinician catch:} Mandatory display of Grad-CAM alongside the score immediately flags off-lesion attention, triggering an image-reacquisition prompt and clinical override.
  \item \textbf{Subgroup Hidden Stratification (Extreme Age or Acral/Head-Neck Sites):} Elevated false-negative rates in underrepresented anatomical sites (palms, soles, scalp, ear) or elderly cohorts ($\ge 75$~years)~\cite{OakdenRayner2020hidden}. \emph{Patient harm:} Moderate-to-High—systematic care disparities. \emph{Clinician catch:} Automated UI caution banner whenever patient metadata fall in a high-variance subgroup, paired with site-specific dermoscopic algorithms (e.g., parallel ridge rule on acral skin).
\end{enumerate}

\subsection{Regulatory Pathway and Approval Obstacles}
Under U.S.\ FDA medical device statutes, an AI triage classifier recommending specialist referral for suspected skin cancer is regulated as Class~II (moderate-to-high risk) \emph{Software as a Medical Device (SaMD)}~\cite{Food2021artificial,Gerke2020need}. Because autonomous skin-cancer diagnosis carries severe false-negative mortality risk, claiming standalone autonomous diagnosis would likely trigger a Class~III \textbf{Premarket Approval (PMA)} pathway requiring multi-year randomized or large-scale prospective confirmatory clinical trials costing tens of millions of dollars. Positioning the product strictly as a \textbf{human-in-the-loop primary-care triage aid} enables either the \textbf{De~Novo} pathway (if no sufficiently similar dermatology AI predicate exists for the exact intended-use population) or a \textbf{510(k)} substantial-equivalence submission referencing cleared dermatology CADe/CADt predicates~\cite{Zhu2022landscape}. Operating as a Laboratory-Developed Test (LDT) is inapplicable to distributed clinical imaging software. Key regulatory obstacles include: (a)~prospective multi-center validation across diverse U.S./global clinics and all six Fitzpatrick skin phototypes; (b)~human-factors and reader-study evidence proving the AI reduces missed cancers without inducing clinician deskilling or excessive biopsy rates; (c)~full design-history documentation under ISO~13485 and IEC~62304; and (d)~submission of a formal \textbf{Predetermined Change Control Plan (PCCP)} defining safe boundaries for post-market retraining~\cite{Food2021artificial}.

\subsection{Privacy, HIPAA Compliance, and Data Governance}
Dermatoscopic images paired with clinical records constitute Protected Health Information (PHI) under the U.S.\ Health Insurance Portability and Accountability Act (HIPAA). Although dermoscopy focuses on localized skin patches rather than full portraits, clinical images can still carry re-identification risks via tattoos, facial landmarks (scalp/ear/face lesions), unique birthmarks, burned-in ruler timestamps, or DICOM headers (PatientName, PatientID, AccessionNumber, StudyDate, InstitutionName, and vendor private elements). Deployment therefore requires end-to-end encryption in transit (mutual TLS~1.3) and at rest (AES-256), role-based access control (RBAC) with immutable audit trails logging every prediction request, Business Associate Agreements (BAAs) with any cloud infrastructure provider, and strict patient informed consent covering secondary model retraining.

At the ingestion boundary, our \texttt{pydicom 3.0.2} adapter implements HIPAA Safe Harbor metadata hygiene: parsing strictly through a minimal non-PHI technical allowlist, stripping private elements, and rejecting files that retain direct identifier tags. Critically, following DICOM PS3.15 Annex~E, we explicitly recognize that header scrubbing alone does \emph{not} constitute complete HIPAA Safe Harbor compliance; production pipelines must add automated optical character recognition (OCR) to detect burned-in pixel text, facial-region screening, and legal privacy-officer verification before any image leaves a covered entity's security perimeter.

\subsection{Fairness, Bias, and Hidden Stratification Audit}
Aggregate performance masks clinically dangerous disparities across patient subgroups~\cite{OakdenRayner2020hidden,Gianfrancesco2018potential,Obermeyer2019dissecting,Rajkomar2018scalable}. Table~\ref{tab:fairness} and Figure~\ref{fig:fairness} present our disaggregated test-set fairness audit across all four clinical metadata attributes available in HAM10000 (\texttt{age\_group}, \texttt{sex}, \texttt{localization\_group}, and \texttt{dx\_type}), reporting sensitivity, specificity, AUC-ROC, and flagging subgroups exhibiting a $>10$~percentage-point drop relative to overall performance.

\begin{table}[htbp]
\centering
\caption{\textbf{Subgroup fairness and hidden stratification audit on the HAM10000 test split.} Alerts flag subgroups with sensitivity or AUC dropping $>10$~pp below the aggregate test score~\cite{OakdenRayner2020hidden}.}
\label{tab:fairness}
\small
\begin{tabular}{llrrrrl}
\toprule
\textbf{Attribute} & \textbf{Subgroup} & \textbf{N} & \textbf{Sens.} & \textbf{Spec.} & \textbf{AUC} & \textbf{Status} \\
\midrule
\multicolumn{7}{c}{\textit{Pending ROCm evaluation run.}} \\
\bottomrule
\end{tabular}

\end{table}

\begin{figure}[htbp]
\centering
\safeincludegraphics[0.96\textwidth]{../figures/fairness_subgroups.png}{Disaggregated sensitivity and specificity across patient age, biological sex, anatomical localization, and diagnostic ground-truth verification method.}
\caption{\textbf{Subgroup fairness and hidden stratification audit across HAM10000 metadata slices.} Exclamation markers highlight subgroups with $>10$~percentage-point degradation relative to aggregate performance.}
\label{fig:fairness}
\end{figure}

Two structural bias findings merit explicit disclosure. First, stratifying by \texttt{dx\_type} reveals \emph{annotation-quality confounding}: biopsy-verified (\texttt{histo}) lesions represent harder, clinically suspicious referrals with higher malignancy prevalence than digital follow-up (\texttt{follow\_up}) nevi, so aggregate metrics blend easy follow-up controls with challenging histopathology cases. Second—and most consequentially for global health equity—HAM10000 originates primarily from Vienna and Queensland clinics with predominantly light-skinned patients (Fitzpatrick phototypes I--III) and does not record Fitzpatrick skin tone~\cite{Tschandl2018dataset,Daneshjou2022disparities}. Deploying this model without targeted data collection and external validation on Fitzpatrick IV--VI populations (such as the Diverse Dermatology Images benchmark~\cite{Daneshjou2022disparities}) would violate algorithmic fairness principles~\cite{Selbst2019fairness} and systematically under-diagnose acral lentiginous melanoma in darker-skinned patients. Concrete mitigations include mandatory multi-ethnic external prospective cohorts, subgroup-calibrated operating thresholds where ethically justified, and explicit labeling of validated phototype scope in FDA Instructions for Use.

\subsection{Explainability, Actionability, and Clinician Trust}
Saliency maps are necessary for error auditing—especially catching border or ruler shortcuts—yet current post-hoc attribution methods exhibit well-documented clinical limitations~\cite{Ghassemi2021false}. Grad-CAM highlights \emph{where} the network looked, not \emph{why} it distinguished melanoma from seborrheic keratosis (e.g., whether it detected atypical pigment networks, blue-white veil, or regression structures). To make explanations genuinely actionable for clinicians~\cite{Samek2021explaining,Caruana2015intelligible}, clinical UI integration must pair localized heatmaps with (1)~calibrated risk estimates and uncertainty intervals, (2)~the top-3 differential diagnoses rather than a single forced label, (3)~retrieval of visually similar biopsy-proven training prototypes, and (4)~clear warning badges whenever the Workshop~4 \texttt{MedConvVAE2D} anomaly detector flags an image as out-of-distribution relative to training dermoscopy.

\subsection{Operational Challenges, Production Monitoring, Drift Detection, and ML Security}
Translating a validated checkpoint into hospital operations demands rigorous MLOps infrastructure and security controls (Figure~\ref{fig:monitoring_drift}). Clinical imaging environments suffer continuous covariate and concept drift from new dermatoscope optics, LED illumination aging, software compression updates, seasonal sun-exposure variation, and shifting referral demographics~\cite{Finlayson2021clinician}. To verify monitoring sensitivity empirically, \texttt{src/evaluate\_medical.py} simulated three realistic hospital shifts on the test cohort (warm illuminant shift, optical defocus $\sigma=2.4$, and $-32\%$ under-exposure with sensor noise) and evaluated our Population Stability Index (PSI) and two-sample Kolmogorov--Smirnov (KS) drift monitors (Table~\ref{tab:drift}).

\begin{table}[htbp]
\centering
\caption{\textbf{Simulated clinical deployment drift monitoring (PSI and KS tests) and automated retraining triggers.}}
\label{tab:drift}
\small
\begin{tabular}{lrrrrl}
\toprule
\textbf{Clinical Shift Scenario} & \textbf{Pred PSI} & \textbf{KS Stat} & \textbf{Mal. AUC} & \textbf{Sens.} & \textbf{Trigger} \\
\midrule
In-Distribution Control (No Shift) & 0.010 & 0.032 & 0.940 & 91.6\% & OK \\
New Dermatoscope Color/Illumination Shift & 0.164 & 0.195 & 0.873 & 55.4\% & INVESTIGATE \\
Optical Defocus / Motion Blur (sigma=2.4) & 0.485 & 0.175 & 0.819 & 56.6\% & RETRAIN\_ALERT \\
Under-Exposure + Sensor Noise (-32\% gain) & 1.733 & 0.313 & 0.578 & 15.7\% & RETRAIN\_ALERT \\
\bottomrule
\end{tabular}

\end{table}

\begin{figure}[htbp]
\centering
\resizebox{0.96\textwidth}{!}{%
\begin{tikzpicture}[
  node distance=0.6cm and 0.8cm,
  blk/.style={rectangle, rounded corners=3pt, draw=OkabeBlue, fill=BoxBg, thick, text width=2.85cm, align=center, minimum height=1.15cm, font=\footnotesize},
  alert/.style={rectangle, rounded corners=3pt, draw=OkabeOrange, fill=orange!8, thick, text width=2.85cm, align=center, minimum height=1.15cm, font=\footnotesize},
  gov/.style={rectangle, rounded corners=3pt, draw=OkabeGreen, fill=teal!8, thick, text width=2.95cm, align=center, minimum height=1.15cm, font=\footnotesize},
  arrow/.style={-{Stealth[length=2.2mm]}, thick, color=OkabeGray}
]
\node[blk] (serve) {\textbf{TLS 1.3 Serving API}\\ Input Validation\\ Latency \& Error Logs};
\node[blk, right=of serve] (log) {\textbf{De-Identified Store}\\ Prediction Probabilities\\ Image Luminance + OOD};
\node[alert, right=of log] (drift) {\textbf{Drift Engine}\\ PSI ($>0.20$) + KS ($p<0.05$)\\ Subgroup FNR Tracking};
\node[gov, right=of drift] (pccp) {\textbf{FDA PCCP Gate}\\ Biopsy Ground Truth\\ Shadow Re-Validation};

\draw[arrow] (serve) -- (log);
\draw[arrow] (log) -- (drift);
\draw[arrow] (drift) -- node[above, font=\scriptsize]{Alert} (pccp);
\draw[arrow] (pccp.south) -- ++(0,-0.65) -| node[pos=0.25, below, font=\scriptsize]{Approved Model Rollout + Signed SHA-256 Artifact} (serve.south);
\end{tikzpicture}
}
\caption{\textbf{Closed-loop clinical monitoring and post-market governance architecture.} Real-time inference logs feed statistical drift detectors (PSI and Kolmogorov--Smirnov tests) and histopathology feedback collectors; retraining proceeds strictly through an FDA Predetermined Change Control Plan (PCCP) validation gate rather than uncontrolled online updates.}
\label{fig:monitoring_drift}
\end{figure}


As designed in Figure~\ref{fig:monitoring_drift}, production serving enforces: (1)~\emph{Input validation and observability} via Prometheus latency/error metrics and structured prediction logs tracking image luminance, sharpness, OOD VAE scores, and predicted probability distributions; (2)~\emph{Automated drift alerting} triggering clinical investigation when $\mathrm{PSI} \in [0.10, 0.20]$ or KS $p < 0.05$, and a formal retraining review when $\mathrm{PSI} > 0.20$ or rolling biopsy-confirmed sensitivity drops $>5$~pp below baseline; (3)~\emph{Controlled retraining under FDA PCCP} using delayed histopathology ground truth collected from EHR pathology reports, validated in shadow mode across all demographic slices before signed artifact promotion; (4)~\emph{ML security hardening} mitigating OWASP ML01 (adversarial perturbations~\cite{Finlayson2019adversarial}), ML06/ML07 (supply-chain and transfer-learning backdoors via SHA-256 signed \texttt{weights\_only=True} checkpoints and dependency auditing), and API rate limiting; and (5)~\emph{Organizational training} educating clinicians on automation bias, subgroup limitations, and explicit override protocols.
